\chapter{User Research}\label{chapter06}

\section*{Introducción breve}

Con el objetivo de comprender las necesidades y problemáticas de potenciales usuarios vinculados al comercio electrónico, se realizó una instancia inicial de User Research orientada a emprendedores, revendedores digitales y usuarios interesados en la detección de tendencias comerciales.

La investigación tuvo como propósito identificar dificultades relacionadas con la búsqueda de productos emergentes, acceso a información histórica y utilización de herramientas de análisis dentro de plataformas ecommerce.

\section{Objetivo de la encuesta}

La encuesta fue diseñada para identificar hábitos de búsqueda de productos, dificultades para detectar tendencias y nivel de utilización de herramientas de análisis comercial.

El relevamiento se realizó mediante una encuesta digital elaborada en Google Forms entre el 26 de mayo y el 16 de julio de 2026 mediante un formulario combinado con preguntas cerradas, escalas de valoración de 1 a 5 y preguntas abiertas breves, permitiendo relevar tanto datos cuantitativos como comentarios cualitativos. El canal de difusión utilizado fue WhatsApp y además contactos personales. El criterio de selección fue no probabilístico por conveniencia, orientado a personas con interés en comercio electrónico, venta digital o compra de productos en marketplaces. Los gráficos completos de resultados generados por Google Forms se incluyen en el Anexo de Resultados de la Encuesta.

\subsection*{Fundamentación de la muestra}

Se obtuvieron 120 respuestas válidas. Para el análisis, las respuestas se segmentaron entre vendedores, compradores y participantes sin rol activo declarado. Las preguntas dirigidas a vendedores se utilizaron para validar la problemática principal de detección tardía de productos, mientras que las preguntas dirigidas a compradores se analizaron como señales complementarias de demanda, precio e influencia de redes sociales.

Como limitación metodológica, la muestra no es representativa de todo el universo de usuarios de Mercado Libre Argentina, ya que se concentra principalmente en CABA y Buenos Aires. Por ese motivo, los resultados se interpretan como exploratorios y útiles para orientar el diseño inicial de la plataforma.

\section{Composición de la muestra}

La primera pregunta permitió dividir a los participantes según su rol principal en Mercado Libre. Sobre 120 respuestas, el 37,5\% indicó ser vendedor, el 50,8\% comprador y el 11,7\% ninguna de las opciones. Esto equivale a 45 vendedores, 61 compradores y 14 participantes sin rol activo declarado.

\begin{table}[ht]
	\centering
	\caption{Composición de la muestra}
	\label{tab:composicion}
	\begin{tabularx}{0.7\textwidth}{@{}X c c@{}}
		\toprule
		\textbf{Rol declarado} & \textbf{Porcentaje} & \textbf{Cantidad aproximada} \\
		\midrule
		Vendedor & 37,5\% & 45 respuestas \\
		Comprador & 50,8\% & 61 respuestas \\
		Ninguna & 11,7\% & 14 respuestas \\
		Total & 100\% & 120 respuestas \\
		\bottomrule
	\end{tabularx}
\end{table}

Las preguntas específicas para vendedores registraron entre 43 y 45 respuestas, debido a que una de las preguntas del bloque vendedor registró dos respuestas menos. Las preguntas específicas para compradores registraron 61 respuestas.

\section{Resultados del segmento vendedor}

El bloque de vendedores permite observar el problema desde el punto de vista de quienes podrían utilizar directamente la plataforma propuesta para tomar decisiones de stock, detectar oportunidades comerciales y analizar tendencias antes de que se consoliden.

\begin{table}[ht]
	\centering
	\caption{Resultados del segmento vendedor}
	\label{tab:resultados-vendedor}
	\begin{tabularx}{\textwidth}{@{}>{\raggedright\arraybackslash}p{3.8cm} >{\raggedright\arraybackslash}p{5.5cm} X@{}}
		\toprule
		\textbf{Pregunta} & \textbf{Resultado principal} & \textbf{Interpretación} \\
		\midrule
		Dificultad para encontrar productos que puedan volverse tendencia & Sobre 43 respuestas, el 69,7\% eligió valores 4 o 5. El promedio fue 3,9 sobre 5. & Existe una dificultad alta para anticipar productos con potencial de tendencia. \\
		Método actual de detección & Tendencias dentro de Mercado Libre: 31,1\%; publicidad: 22,2\%; recomendaciones: 20\%; redes sociales: 13,3\%; otros: 13,3\%. & Los vendedores combinan varias fuentes, pero no aparece un método único ni sistemático. \\
		Percepción de detección tardía & 40\% respondió Sí, 48,9\% A veces y 11,1\% No. & La mayoría considera que las tendencias se detectan tarde o, al menos, no siempre a tiempo. \\
		Información más útil en una herramienta de análisis & Predicción de crecimiento: 28,9\%; competencia: 26,7\%; categorías de crecimiento: 20\%; historial de precios: 13,3\%; productos emergentes: 11,1\%. & La necesidad principal no es solo ver productos, sino anticipar crecimiento y analizar competencia. \\
		Disposición a usar una herramienta automatizada & 68,9\% respondió Sí, 26,7\% Tal vez y 4,4\% No. & Existe una aceptación alta de una solución automatizada para detectar productos emergentes. \\
		\bottomrule
	\end{tabularx}
\end{table}

\section{Resultados del segmento comprador}

El bloque comprador permite analizar señales de comportamiento del lado de la demanda. Aunque el usuario principal de la plataforma puede ser el vendedor o emprendedor, entender qué motiva a los compradores ayuda a definir qué variables y visualizaciones pueden resultar relevantes.

\begin{table}[ht]
	\centering
	\caption{Resultados del segmento comprador}
	\label{tab:resultados-comprador}
	\begin{tabularx}{\textwidth}{@{}>{\raggedright\arraybackslash}p{3.8cm} >{\raggedright\arraybackslash}p{5.5cm} X@{}}
		\toprule
		\textbf{Pregunta} & \textbf{Resultado principal} & \textbf{Interpretación} \\
		\midrule
		Motivación para comprar productos nuevos o novedosos & Recomendaciones: 26,2\%; redes sociales: 21,3\%; novedad: 19,7\%; precio: 16,4\%; publicidad: 16,4\%. & Las señales sociales (recomendaciones y redes) desplazan al precio como factor principal en la compra de productos novedosos, hallazgo relevante para las variables del modelo. \\
		Percepción de resultados repetidos en Mercado Libre & El 36,1\% eligió 5, el 23\% eligió 4, el 24,6\% eligió 3, el 14,8\% eligió 2 y el 1,5\% eligió 1. Promedio aproximado: 3,8 sobre 5. & Existe una percepción alta de repetición o concentración de resultados destacados (el 59,1\% eligió valores 4 o 5). \\
		Influencia de TikTok e Instagram & Promedio aproximado: 3,3 sobre 5. El 47,5\% eligió valores 4 o 5. & La influencia de redes sociales es moderada, pero relevante como señal complementaria para entender demanda. \\
		\bottomrule
	\end{tabularx}
\end{table}

\section{Resultados generales}

\begin{table}[ht]
	\centering
	\caption{Resultados generales}
	\label{tab:resultados-generales}
	\begin{tabularx}{\textwidth}{@{}>{\raggedright\arraybackslash}p{2.5cm} >{\raggedright\arraybackslash}p{6cm} X@{}}
		\toprule
		\textbf{Variable} & \textbf{Resultado} & \textbf{Lectura} \\
		\midrule
		Rango de edad & 18 a 30: 43,3\%; 31 a 50: 40,8\%; más de 50: 10,8\%; menos de 18: 5,1\%. & La muestra se concentra principalmente en adultos jóvenes y adultos en edad laboral. \\
		Ubicación & Sobre 119 respuestas, se agruparon variantes escritas manualmente: Buenos Aires / Provincia de Buenos Aires aprox. 40\%; CABA aprox. 28\%; respuestas genéricas (``Argentina'') u otras provincias aprox. 32\%. & La muestra se concentra en CABA y Buenos Aires, lo cual debe considerarse como una limitación del relevamiento. \\
		Comentarios abiertos & Aparecieron menciones a comparación de precios, opiniones del vendedor, recomendación de productos nuevos y necesidad de una herramienta visual. & Los comentarios refuerzan la utilidad de una plataforma visual, orientada a precios, recomendaciones y análisis comercial. \\
		\bottomrule
	\end{tabularx}
\end{table}

\section{Principales datos obtenidos}

\begin{itemize}
	\item Existe una dificultad alta para anticipar productos tendencia: el 69,7\% de los vendedores que respondieron esa pregunta seleccionó valores 4 o 5.
	\item La detección tardía aparece como un problema relevante: el 40\% respondió que las tendencias muchas veces se detectan demasiado tarde y el 48,9\% respondió que eso ocurre a veces (88,9\% combinado).
	\item La aceptación de una herramienta automatizada es alta: el 68,9\% de los vendedores indicó que la usaría y un 26,7\% adicional respondió ``tal vez'' (solo el 4,4\% la descartó).
	\item Los vendedores valoran especialmente la predicción de crecimiento y el análisis de competencia, lo que refuerza la necesidad de una plataforma que no se limite a mostrar rankings actuales.
	\item Desde el lado comprador, las recomendaciones (26,2\%) y las redes sociales (21,3\%) aparecen como los factores principales de decisión, seguidas por la novedad (19,7\%) y el precio (16,4\%); esto refuerza el valor de las señales sociales como variables complementarias para interpretar oportunidades comerciales.
	\item La muestra se concentra en CABA y Buenos Aires, por lo que los resultados deben entenderse como exploratorios y no generalizables a todo el país.
\end{itemize}

Los resultados obtenidos respaldan la problemática inicial del proyecto. Los vendedores manifiestan dificultad para anticipar tendencias, perciben que muchas oportunidades se detectan tarde y muestran alta disposición a utilizar una herramienta automatizada. Al mismo tiempo, los compradores confirman la influencia de variables como recomendaciones, redes sociales y precio. En conjunto, estos hallazgos justifican el desarrollo de una plataforma orientada a detectar productos emergentes mediante históricos, métricas temporales y visualización analítica.

\section{Entrevistas en profundidad}

Como complemento cualitativo de la encuesta, se diseñaron tres entrevistas semiestructuradas a usuarios reales de Mercado Libre: dos vendedores y un comprador, seleccionados por muestreo intencional entre personas con actividad verificable en la plataforma (no se buscan perfiles profesionales ni expertos, sino usuarios comunes, en línea con el público objetivo de la herramienta). El propósito es profundizar en el porqué de los patrones cuantitativos observados en la encuesta: la detección tardía de tendencias (88,9\% combinado), la alta disposición a usar una herramienta automatizada (68,9\%) y el peso de las señales sociales en la decisión de compra (47,5\%). Cada entrevista dura entre 20 y 30 minutos, se registra con consentimiento informado y se analiza mediante síntesis temática. Los hallazgos se integrarán a esta sección una vez realizadas las entrevistas. El diseño completo de las tres guías de entrevista (objetivos y las 10 preguntas de cada una) se incluye en el Anexo de Entrevistas.

\section{User persona}

A partir de los hallazgos reales de la encuesta se construyen dos user personas preliminares. No representan personas individuales encuestadas, sino perfiles sintéticos elaborados a partir de patrones observados en las respuestas. Su objetivo es conectar los resultados del relevamiento con necesidades concretas de diseño de la plataforma.

\subsection*{User Persona 1: Martín, vendedor que necesita anticiparse}

\begin{table}[ht]
	\centering
	\caption{User Persona 1 -- Martín}
	\label{tab:persona-martin}
	\begin{tabularx}{\textwidth}{@{}>{\raggedright\arraybackslash}p{3cm} X@{}}
		\toprule
		\textbf{Campo} & \textbf{Descripción} \\
		\midrule
		Perfil & Martín tiene 29 años, vive en Buenos Aires y vende productos por Mercado Libre. Representa al segmento vendedor de la encuesta, que declaró dificultades altas para encontrar productos que puedan volverse tendencia. Tiene experiencia básica/intermedia en ventas online y consulta distintas fuentes antes de decidir qué productos incorporar. \\
		Necesidad & Necesita detectar productos con potencial de crecimiento antes de que se consoliden y antes de que aumente demasiado la competencia. También necesita comparar la evolución del producto con su categoría para evitar decisiones basadas solo en una aparición aislada. \\
		Frustración & Siente que muchas veces llega tarde a las tendencias: cuando identifica un producto atractivo, ya existen muchos vendedores ofreciendo lo mismo y los márgenes se reducen. \\
		Objetivo & Tomar decisiones de stock con más información, priorizando productos que muestren crecimiento sostenido y señales tempranas de demanda. \\
		Relación con la plataforma propuesta & Sería un usuario directo del dashboard. Consultaría productos emergentes, predicción de crecimiento, evolución de precios y nivel de competencia para decidir qué vender o qué analizar con mayor profundidad. La plataforma le permitiría priorizar oportunidades según señales históricas, no solamente según intuición o publicaciones virales. \\
		\bottomrule
	\end{tabularx}
\end{table}

\subsection*{User Persona 2: Luciana, compradora digital con comportamiento influido por precio y redes}

\begin{table}[ht]
	\centering
	\caption{User Persona 2 -- Luciana}
	\label{tab:persona-luciana}
	\begin{tabularx}{\textwidth}{@{}>{\raggedright\arraybackslash}p{3cm} X@{}}
		\toprule
		\textbf{Campo} & \textbf{Descripción} \\
		\midrule
		Perfil & Luciana tiene 32 años, vive en CABA y compra habitualmente en Mercado Libre. Representa al segmento comprador de la encuesta, donde las recomendaciones y las redes sociales aparecen como los factores más relevantes para decidir compras de productos novedosos. Su comportamiento representa señales indirectas de demanda que los vendedores pueden observar para interpretar interés de mercado. \\
		Necesidad & Quiere encontrar productos nuevos o convenientes sin depender únicamente de resultados destacados repetidos o recomendaciones superficiales. Necesita contar con información clara que le permita distinguir entre productos realmente convenientes y productos destacados solo por visibilidad comercial. \\
		Frustración & Percibe que muchas veces aparecen los mismos productos destacados y que no siempre es fácil comparar alternativas, precios u opiniones del vendedor. \\
		Objetivo & Comprar productos novedosos o recomendados con mayor confianza, comparando precio, reputación y señales de interés real. \\
		Relación con la plataforma propuesta & Aunque no sería necesariamente la usuaria principal del sistema, su comportamiento ayuda a definir qué señales de demanda deberían observar los vendedores: precio, redes sociales, recomendaciones y aparición repetida de productos. Sus hábitos permiten definir variables complementarias para el sistema, especialmente precio, repetición de publicaciones, reputación y presencia en redes sociales. \\
		\bottomrule
	\end{tabularx}
\end{table}

\section{Diagrama de Gantt}

El cronograma de actividades del proyecto se presenta en el Anexo de Cronograma de Actividades.

\section{Gestión de datasets}

El dataset del proyecto será construido de manera propia a partir de consultas periódicas a la API oficial de Mercado Libre Argentina. No se utilizará scraping, fuentes no autorizadas ni datos privados de usuarios. Los datos crudos provendrán de información pública disponible mediante endpoints oficiales, principalmente publicaciones, categorías, precios y atributos asociados a productos.

En esta etapa inicial, el volumen del dataset se plantea como estimado, ya que el histórico comenzará a construirse durante el desarrollo del prototipo. Primero se trabajará con una muestra controlada de productos y categorías, y luego podrá ampliarse incorporando mayor frecuencia de captura y más registros temporales.

La confiabilidad del origen se basa en el uso de la API oficial de Mercado Libre. Los datos serán almacenados en PostgreSQL para conservar registros estructurados y permitir comparaciones entre distintos períodos.

El tratamiento preliminar incluirá control de duplicados, validación de campos obligatorios, revisión de valores nulos, normalización de categorías y registro de fecha y hora de captura. De manera preliminar, se considerarán variables como identificador de publicación, título, categoría, precio, fecha de captura, frecuencia de aparición y cantidad de publicaciones asociadas.

Como limitación inicial, el proyecto no contará con históricos previos provistos por Mercado Libre. Por eso, el valor del dataset aumentará a medida que el sistema acumule registros en el tiempo.
